% Synthetic thesis: all data, models, and citation keys are fictional test material.
\documentclass{ctexbook}
\begin{document}
\chapter{方法与实验}
本章介绍合成样本的处理流程，并记录用于核对润色范围的占位实验条件。

\section{局部误差分析}
本节比较两个占位模型的预测误差，所有结论仅限于本节所述的合成样本。

\subsection{误差比较}
本文不试图说明其他装置的运行效果。值得注意的是，本研究在相同条件下比较了模型 A 和模型 B 的预测误差，共统计 20 个样本。模型 A 的 MAE 为 0.12，模型 B 的 MAE 为 0.15，测量间隔为 5\,s。该结果可能支持模型 A 在当前样本上的误差优势\cite{demo2026a,demo2026b}。

误差按式\eqref{eq:demo-error}计算，输入与输出的对应关系见图\ref{fig:demo-flow}。变量 $e=y-\hat{y}$ 表示单个样本的残差。
\begin{equation}
e=y-\hat{y}
\label{eq:demo-error}
\end{equation}

\subsection{邻域说明}
后续分析只记录占位样本的排列规则并保留独立的讨论范围。
本小节不仅记录样本顺序，而且列出相应索引，但没有据此推断其他场景的结果。

\chapter{过程记录}
本章使用另一组占位描述展示方法段的语句结构，不增加前章结论的适用范围。

\section{处理步骤}
本节记录数据进入处理流程后的操作顺序，供单元润色时保留方法步骤。

\subsection{样本处理}
首先读取占位容器中的记录，再按时间排序并识别空缺位置，然后统计各组样本的残差并计算组内均值，最后把组号和计算结果写入索引以便复核。
\end{document}
